\subsection{Canonical Schur covers and their relations}
\label{sec:schur-presentation}

For the remainder of this section $K$ has characteristic zero, and
$B=K\oplus\mathfrak m$ is any augmented commutative algebra. Put
$S_B=\Sym(\mathfrak m)$; its generator $T_b$, linear in $b\in\mathfrak m$,
acts on $P_n(B)=B^{\otimes n}$ by $p_n(b)=\sum_jb^{(j)}$.
The trace presentation in Section~\ref{sec:selected-quotients} identifies
$Q_n(B)$ with $P_n(B)/(S_B)_+P_n(B)$.

The augmentation decomposes $P_n(B)$ by occupied positions:
\[
 P_n(B)=\bigoplus_{j=0}^nP_n^{\langle j\rangle}(B),\qquad
 P_n^{\langle j\rangle}(B)=
       \bigoplus_{\substack{I\subseteq[n]\\|I|=j}}\mathfrak m^{\otimes I},
\]
with units elsewhere. If $e_\mu$ is the central label idempotent, define
the actual subspace
\[
 L_n(B)=\bigoplus_{\mu\vdash n}
          e_\mu P_n^{\langle n-\mu_1\rangle}(B)\subseteq P_n(B).
\]

\begin{lemma}[Minimum occupancy]
\label{lem:minimum-occupancy-cover}
The quotient map $L_n(B)\to Q_n(B)$ is surjective. For
$\mu=(t,\tau)\vdash n$, its label-$\mu$ source multiplicity is
$\Schur_\tau\mathfrak m$.
\end{lemma}
\begin{proof}
Put $s=|\tau|$. Selected elimination makes
$B^{\otimes s}\to Q_n(B)^{S_t}$ surjective. A selected tensor with
a unit occupies at most $s-1$ positions, so its global orbit is induced
from a subgroup containing $S_{t+1}$. Its label first row is therefore
at least $t+1$, and its $\mu$ projection is zero. Fully occupied
tensors cover the $S_t$-fixed part of label $\mu$, and its $S_n$-span
covers the whole isotype. For the trivial label this says
$Q_n(B)^{S_n}=K$.

Finally,
$P_n^{\langle s\rangle}=\operatorname{Ind}_{S_s\times S_t}^{S_n}
(\mathfrak m^{\otimes s}\otimes\mathbf1)$.
If $|\nu|=s$ and $\mu/\nu$ is horizontal, Pieri interlacing gives
$\nu_i\ge\tau_i$; equality of sizes forces $\nu=\tau$.
\end{proof}

For $\mu=(t,\tau)$ write $a=\tau_1$, with $a=0$ for the empty tail,
and put
\[
 \tau^+=(t+1,\tau_2,\tau_3,\ldots),\qquad
 d_1=|\tau^+|=n+1-a,
 \qquad M_\mu(B)=\Hom_{S_n}(S^\mu,Q_n(B)).
\]
Fix the usual one-dimensional Pieri identification lines.

\begin{theorem}[The two-Schur presentation]
\label{thm:two-schur-presentation}
There are natural exact presentations
\begin{align*}
 S_B\otimes\Schur_{\tau^+}\mathfrak m
 &\xrightarrow{R_{\mu,B}}
 S_B\otimes\Schur_\tau\mathfrak m
 \xrightarrow{\Phi_{\mu,B}}\Schur_\mu(B)\longrightarrow0,\\
 \Schur_{\tau^+}\mathfrak m
 &\xrightarrow{\beta_{\mu,B}}\Schur_\tau\mathfrak m
               \longrightarrow M_\mu(B)\longrightarrow0.
\end{align*}
Here $\Phi$ is actual multiplication by traces on minimum occupancy.
All maps commute with arbitrary augmented algebra maps, including
ungraded and noninjective maps, and preserve compatible gradings and
auxiliary symmetries. The left-hand maps need not be injective.
\end{theorem}

\begin{proof}
Filter $F_\mu=S_B\otimes\Schur_\tau\mathfrak m$ by input count:
$j$ formal traces and $|\tau|$ inputs have count $j+|\tau|$.
Filter the target by occupied positions. Collisions lower this count,
so the associated-graded map is the same cover for the square-zero
algebra $B_0=K\oplus\mathfrak m$. Giving $\mathfrak m$ degree one,
Lemma~\ref{lem:minimum-occupancy-cover} makes this cover surjective by
degree induction. Pieri and branching off the constant factor give
\begin{align*}
 F_\mu&=\bigoplus_{\nu/\tau\ \mathrm{horizontal}}\Schur_\nu\mathfrak m,\\
 \Schur_\mu(K\oplus\mathfrak m)
   &=\bigoplus_{\substack{\nu/\tau\ \mathrm{horizontal}\\\nu_1\le t}}
                                      \Schur_\nu\mathfrak m.
\end{align*}
Let $G_\mu$ be the sum of the retained source summands $\nu_1\le t$,
and $Z_\mu$ the omitted sum. These are canonical Schur isotypic
subfunctors. The square-zero map is zero on $Z_\mu$ and an isomorphism
on $G_\mu$. For $B$, the restriction
$\Psi_B=\Phi_{\mu,B}|_{G_\mu}$ has the same diagonal and only
lower-count terms. Retained counts lie between $|\tau|$ and $n$,
so its triangular inverse is finite. The entire kernel is the graph
\begin{equation}\label{eq:schur-relation-graph}
 Z_\mu\xrightarrow{\sim}\ker\Phi_{\mu,B},\qquad
 z\longmapsto z-\Psi_B^{-1}\Phi_{\mu,B}(z).
\end{equation}

Every omitted $\nu$ contains $\tau^+$ and $\nu/\tau^+$ is horizontal.
The symmetric Pieri composition lemma \cite[Lemma~2.1]{SamWeyman}
and multiplicity one therefore give
$Z_\mu=S_B\cdot\Schur_{\tau^+}\mathfrak m$ for the square-zero map.
Restrict \eqref{eq:schur-relation-graph} to this generating summand
and extend $S_B$-linearly. Subtracting lifted generators with a prescribed
leading term lowers count, so finite induction proves that they generate
the actual kernel. Reduction modulo $(S_B)_+$ gives the second
presentation, with the convention
$\beta_{\mu,B}=-\operatorname{ev}_{T=0}R_{\mu,B}$.

Schur projections, multiplication, and the inverse of the natural
isomorphism $\Psi_B$ commute with augmented maps. For infinite-dimensional
$\mathfrak m$, the square-zero constructions pass through finite-dimensional
vector subspaces; the formulas use finitely many products on each input.
\end{proof}

The graph is an identification of filtered vector spaces, not generally
of $S_B$-modules. For example,
\[
 M_{(2,2)}(B)=\Sym^2\mathfrak m\big/
 \Span\{ab\odot c+ac\odot b+bc\odot a:a,b,c\in\mathfrak m\}.
\]
The standard label gives $\mathfrak m/\mathfrak m^n$, as above.
If $\mathfrak m^2=0$, every $\beta$ is zero and
$M_{(t,\tau)}(B)=\Schur_\tau\mathfrak m$ in every rank.

\begin{proposition}[The first obstruction]
\label{prop:first-schur-obstruction}
Suppose $B$ is connected positively graded. The actual free cover on
label $\mu$ is an isomorphism in degrees less than $d_1=n+1-\tau_1$,
and every positive moment Tor group on this label vanishes there.
In degree $d_1$ there is a natural exact sequence
\begin{align*}
 0&\longrightarrow
 \Hom_{S_n}\bigl(S^\mu,\Tor_1^{S_B}(P_n(B),K)_{d_1}\bigr)
 \longrightarrow\Schur_{\tau^+}(B_1)\\
 &\xrightarrow{\beta_{\mu,B}}(\Schur_\tau\mathfrak m)_{d_1}
       \longrightarrow M_\mu(B)_{d_1}\longrightarrow0.
\end{align*}
All higher Tor groups on this label vanish in degree $d_1$.
More generally, $i$th Tor vanishes in degrees less than $d_1+i-1$.
\end{proposition}
\begin{proof}
Every omitted partition has size at least $d_1$, with equality only
for $\tau^+$. The graph construction identifies the first possible
kernel piece with $\Schur_{\tau^+}(B_1)$. Tensor the free-cover
sequence with $K$. Positive coefficients cannot reach that degree,
giving the exact sequence. The Koszul chains of the kernel in
homological degree $i-1$ have degree at least $d_1+i-1$.
\end{proof}

\subsection{A full resolution and its polynomial first degrees}
\label{sec:schur-resolution-statements}

The presentation extends in every homological degree. Extend the parts
of $\mu$ by zero and set
\begin{equation}\label{eq:schur-resolution-shapes}
 \rho_i=(\mu_1+1,\ldots,\mu_i+1,\mu_{i+2},\mu_{i+3},\ldots),
 \qquad d_i=|\rho_i|=n-\mu_{i+1}+i.
\end{equation}
Thus $\rho_0=\tau$ and $\rho_1=\tau^+$. The operation deletes row
$i+1$ and adds one box to each preceding row.

\begin{theorem}[A natural filtered Schur resolution]
\label{thm:full-schur-resolution}
For every augmented commutative $B$ in characteristic zero,
$\Schur_\mu(B)$ has a natural $S_B$-free resolution with terms
\[
                       F_i=S_B\otimes\Schur_{\rho_i}\mathfrak m.
\]
The augmentation and first differential are those of
Theorem~\ref{thm:two-schur-presentation}. Its associated graded for
input count is the classical pure Pieri resolution; the lifted
resolution need not be minimal. Modulo $(S_B)_+$ it gives a complex
\begin{equation}\label{eq:small-schur-moment-complex}
 \cdots\longrightarrow\Schur_{\rho_2}\mathfrak m
 \xrightarrow{\partial_2}\Schur_{\rho_1}\mathfrak m
 \xrightarrow{\partial_1}\Schur_\tau\mathfrak m
\end{equation}
computing $\Hom_{S_n}(S^\mu,\Tor_i^{S_B}(P_n(B),K))$.
The differential $\partial_i$ uses exactly
$e_i=\mu_i-\mu_{i+1}+1$ multiplications. It is natural for all
augmented algebra maps and preserves compatible gradings.
For $\dim\mathfrak m=q$, the terms with $i>q$ vanish.
\end{theorem}

Appendix~\ref{app:schur-resolution} constructs the maps by a natural
creation--annihilation contraction. Its Laplacian has a
dimension-independent content eigenvalue, and multiplication is a
finite perturbation. This constructs the full resolution without choices
of higher splittings.

\paragraph{Nilpotent augmentation ideals.}
If $\mathfrak m^L=0$ with $L\ge2$, then $\partial_i=0$ whenever
$d_i>(L-1)d_{i-1}$: every term partitions its $d_i$ inputs into
$d_{i-1}$ products, and one product must contain at least $L$ inputs.
In particular, the canonical presentation gives
\[
 M_{(n-|\tau|,\tau)}(B)\cong\Schur_\tau\mathfrak m
 \qquad\bigl(n\ge(L-1)|\tau|+\tau_1\bigr).
\]
This equality holds in every internal degree whenever $B$ is graded.
For square-zero algebras it holds at every rank supporting the label.

\begin{corollary}[Sharp universal homological support]
\label{cor:universal-schur-tor-support}
For connected positively graded $B$ and $i\ge1$, label $\mu$ in
degree-$D$ moment homology requires
\[
 D\ge n+i-\mu_{i+1},\qquad
 D\ge n+i-\left\lfloor\frac n{i+1}\right\rfloor.
\]
Both bounds are attained by square-zero algebras with sufficiently
large degree-one space.
\end{corollary}
\begin{proof}
The $i$th term of \eqref{eq:small-schur-moment-complex} begins in
internal degree $d_i$, and $\mu_{i+1}\le\lfloor n/(i+1)\rfloor$.
For square-zero multiplication all differentials are zero. A partition
maximizing its $(i+1)$st part attains the second threshold.
\end{proof}

For a polynomial algebra, the first possible degree behaves better.
The proof uses an ungraded specialization to orthogonal idempotents;
the full naturality in Theorem~\ref{thm:two-schur-presentation} is essential.

\begin{theorem}[Exactness in the first possible homological degree]
\label{thm:polynomial-schur-fronts}
Let $A=\Sym(V)$ and $S_A=\Sym(A_+)$. If $\mu\ne(n)$ and $i\ge1$,
then
\[
 \Hom_{S_n}\bigl(S^\mu,\Tor_i^{S_A}(A^{\otimes n},K)_D\bigr)=0
                         \qquad(D\le n+i-\mu_{i+1}).
\]
The same conclusion holds for $\mu=(n)$ and $i\ge2$.
The sole exceptional first degree is
\[
 \Hom_{S_n}\bigl(\mathbf1,
       \Tor_1^{S_A}(A^{\otimes n},K)_{n+1}\bigr)=\Sym^{n+1}V.
\]
Equivalently, every supported all-linear restriction
$\partial_i:\Schur_{\rho_i}V\to
(\Schur_{\rho_{i-1}}A_+)_{d_i}$ is injective, except for
$\mu=(n),i=1$, where it is zero.
\end{theorem}

\begin{proof}
Take the split augmented algebra $C=K^{q+1}$, with augmentation its
zeroth coordinate and augmentation ideal spanned by the orthogonal
idempotents $u_1,\ldots,u_q$. On a primitive tensor idempotent of
$C^{\otimes n}$, the moment $p_n(u_a)$ acts by the number of positions
of color $a$. Unless every position has color zero, some such count
is a positive integer, and its Koszul complex contracts. Consequently
\begin{equation}\label{eq:idempotent-moment-homology}
 Q_n(C)=K,\qquad
 \Tor_i^{\Sym(C_+)}(C^{\otimes n},K)
                      =\Lambda^i(C_+)\otimes\mathbf1.
\end{equation}

Suppose $\mu\ne(n)$ and $\Schur_{\rho_i}K^q\ne0$.
The complex \eqref{eq:small-schur-moment-complex} for $C$ is exact
in every degree on this label. Its term dimensions are those of the
square-zero pure resolution. Localizing that resolution at the fraction
field of $\Sym(K^q)$ also gives an exact complex: the square-zero
coefficient module has finite length. Exactness determines the ranks
of all differentials recursively from the term dimensions. Each
supported pure Pieri differential is nonzero, so $\partial_i$ for
$C$ is nonzero as well.

Set $q=\dim V$ and extend a linear isomorphism $V\to C_+$ to an
augmented algebra map $A\to C$. Naturality now makes the restriction
of $\partial_i$ to $\Schur_{\rho_i}V$ nonzero. It is $\GL(V)$-equivariant
with irreducible source, hence injective. If that Schur module is zero,
there is nothing to prove.

For $\mu=(n)$ and $i\ge2$, take $q=i$. The top small-complex term is
$\Schur_{(n+1,1^{i-1})}K^i$, of dimension
$\binom{n+i-1}{n}>1$, whereas its kernel has dimension one by
\eqref{eq:idempotent-moment-homology}. Thus its differential is nonzero.
The same specialization, followed by naturality in the row space,
proves injectivity in every supporting dimension. For $i=1$ the
target is $K$, so the positive-degree map is zero.

Finally, the degree-$d_i$ part of the $i$th Schur term uses only
$V=A_1$, and the next term starts strictly above $d_i$. The assertions
follow from the small complex and the preceding degree bounds.
\end{proof}

\begin{theorem}[The first character correction in every degree]
\label{thm:first-unstable-specht-wall}
For every nonempty $\tau\vdash s$, every $D>s$, and
\[
 n=D+\tau_1-1,\qquad
 \mu=(D-s+\tau_1-1,\tau),\qquad
 \widehat\tau=(D-s+\tau_1,\tau_2,\tau_3,\ldots),
\]
there is a natural exact sequence
\begin{equation}\label{eq:first-unstable-specht-wall}
 0\longrightarrow\Schur_{\widehat\tau}V
 \xrightarrow{\beta_{\mu,A}}(\Schur_\tau A_+)_D
 \longrightarrow\Hom_{S_n}(S^\mu,R_{n,k}[D])\longrightarrow0.
\end{equation}
All positive moment homology on this label vanishes through degree $D$.
Thus, when $\dim V\ge\ell(\tau)$, the stable threshold
$n\ge D+\tau_1$ is sharp in every degree $D>s$.
\end{theorem}
\begin{proof}
Here $D=d_1$ and $\rho_1=\widehat\tau$. Apply
Proposition~\ref{prop:first-schur-obstruction} and
Theorem~\ref{thm:polynomial-schur-fronts}. The correction has length
$\ell(\tau)$ and is nonzero under the stated dimension hypothesis.
\end{proof}

This identifies the canonical relation map at the first rank below
the stable threshold.
